\documentclass[11pt,a4paper]{article}

% ================ Encodage et langue ================
\usepackage[utf8]{inputenc}
\usepackage[T1]{fontenc}
\usepackage[french]{babel}
\usepackage{lmodern}
\usepackage{textcomp}

% ================ Mise en page ================
\usepackage[margin=2.2cm]{geometry}
\usepackage{parskip}
\setlength{\parindent}{0pt}
\setlength{\parskip}{0.6em}

% ================ Couleurs et liens ================
\usepackage[table]{xcolor}
\definecolor{primary}{HTML}{0B5394}
\definecolor{accent}{HTML}{1F77B4}
\definecolor{light}{HTML}{F2F2F2}
\definecolor{codebg}{HTML}{F6F8FA}
\definecolor{rulec}{HTML}{D0D7DE}

\usepackage[colorlinks=true, linkcolor=primary, urlcolor=primary, citecolor=primary]{hyperref}

% ================ Tableaux ================
\usepackage{longtable, booktabs, array, tabularx}
\renewcommand{\arraystretch}{1.25}
\arrayrulecolor{rulec}

% ================ Listes ================
\usepackage{enumitem}
\setlist{itemsep=2pt, topsep=4pt}

% ================ Code et citations ================
\usepackage{fancyvrb}
\usepackage{framed}
\definecolor{shadecolor}{named}{light}

% Encadrés pour les blocs > (quotes)
\usepackage{mdframed}
\newmdenv[
  linecolor=primary,
  linewidth=2pt,
  topline=false, bottomline=false, rightline=false,
  innerleftmargin=10pt, innerrightmargin=8pt,
  innertopmargin=6pt, innerbottommargin=6pt,
  backgroundcolor=codebg,
  skipabove=8pt, skipbelow=8pt
]{quotebox}
\renewenvironment{quote}{\begin{quotebox}}{\end{quotebox}}

% ================ Titres ================
\usepackage{titlesec}
\setcounter{secnumdepth}{0}
\titleformat{\section}{\Large\bfseries\color{primary}}{\thesection.}{0.6em}{}
\titleformat{\subsection}{\large\bfseries\color{accent}}{\thesubsection}{0.6em}{}
\titleformat{\subsubsection}{\normalsize\bfseries}{\thesubsubsection}{0.6em}{}

% ================ En-têtes et pieds de page ================
\usepackage{fancyhdr}
\pagestyle{fancy}
\fancyhf{}
\fancyhead[L]{\small\color{gray}ESPRIT School of Business --- L3 Business Computing BI --- Projet Intégré}
\fancyfoot[C]{\small\thepage}
\renewcommand{\headrulewidth}{0.3pt}
\renewcommand{\footrulewidth}{0pt}

% ================ Pandoc helpers ================
\providecommand{\tightlist}{\setlength{\itemsep}{0pt}\setlength{\parskip}{0pt}}
\providecommand{\passthrough}[1]{#1}

\usepackage{calc}
\providecommand{\real}[1]{#1}


% ================ Métadonnées ================
\title{\color{primary}\textbf{Guide Étudiant --- Démarche et Outils par
Étape}}
\author{Aymen Ben Brik et Hela Hammami}
\date{2026}

\begin{document}

\begin{titlepage}
\centering
\vspace*{1.5cm}

{\Large\bfseries ESPRIT School of Business}\\[0.3cm]
{\large Licence 3 --- Business Computing BI}\\[2cm]

{\color{gray}\large Projet Intégré}\\[0.3cm]
{\color{gray}\large Analyse et Prédiction du Churn Client}\\[2.2cm]

\rule{0.85\linewidth}{1.2pt}\\[0.6cm]
{\Huge\bfseries\color{primary} Guide Étudiant --- Démarche et Outils par
Étape\par}
\rule{0.85\linewidth}{1.2pt}\\[2cm]

{\large Tuteurs du projet}\\[0.2cm]
{\Large\bfseries Aymen Ben Brik et Hela Hammami}\\[1cm]
{\large 2026}

\vfill
{\color{gray}\small Document à destination des étudiants}
\end{titlepage}

\tableofcontents
\newpage

Ce guide vous accompagne \textbf{étape par étape} dans la réalisation du
projet. Pour chaque phase, vous trouverez :

\begin{itemize}
\tightlist
\item
  Le socle \textbf{obligatoire} à réaliser et son échéance.
\item
  Les \textbf{outils imposés} et les \textbf{bonus}, avec une
  proposition d'outils pour le reste (justifiez vos choix dans le
  rapport).
\end{itemize}

\begin{quote}
\textbf{Conseil général} : traitez d'abord tout ce qui est
\textbf{obligatoire}. Les bonus ne rapportent des points que si le socle
est solide.
\end{quote}

\begin{center}\rule{0.5\linewidth}{0.5pt}\end{center}

\section{Étape 0 --- Mise en place de
l'environnement}\label{uxe9tape-0-mise-en-place-de-lenvironnement}

\subsection{Objectifs}\label{objectifs}

\begin{itemize}
\tightlist
\item
  Préparer un environnement de travail reproductible.
\item
  Mettre en place le dépôt GitHub et la collaboration.
\end{itemize}

\subsection{À faire}\label{uxe0-faire}

\begin{itemize}
\tightlist
\item
  Créer un dépôt \textbf{GitHub} sans aucune donnée.
\item
  Inviter les tuteurs en tant que collaborateurs.
\item
  Créer un environnement virtuel Python (\texttt{venv} ou
  \texttt{conda}) et rédiger un \texttt{requirements.txt}.
\item
  Rédiger un premier \texttt{README.md} : description, structure du
  dépôt, instructions d'installation.
\item
  Créer un \texttt{.gitignore} qui exclut \textbf{toutes les données}
  (\texttt{data/}, \texttt{*.csv}, \texttt{*.xlsx}, \texttt{*.pkl},
  \texttt{*.pbix}) et les fichiers d'environnement (\texttt{venv/},
  \texttt{\_\_pycache\_\_/}, \texttt{.env}).
\end{itemize}

\subsection{Outils proposés}\label{outils-proposuxe9s}

{\def\LTcaptype{none} % do not increment counter
\begin{longtable}[]{@{}ll@{}}
\toprule\noalign{}
Catégorie & Options \\
\midrule\noalign{}
\endhead
\bottomrule\noalign{}
\endlastfoot
Versioning & \textbf{Git + GitHub} (obligatoire) \\
IDE & VS Code, PyCharm \\
Notebooks & Jupyter Lab, VS Code Notebooks \\
Gestion de projet & GitHub Projects, Trello, Notion \\
\end{longtable}
}

\begin{center}\rule{0.5\linewidth}{0.5pt}\end{center}

\section{Étape 1 --- Collecte et compréhension des données (semaine du
9/11)}\label{uxe9tape-1-collecte-et-compruxe9hension-des-donnuxe9es-semaine-du-911}

\subsection{Obligatoire}\label{obligatoire}

\begin{itemize}
\tightlist
\item
  Charger les fichiers, vérifier le grain et les clés, croiser avec les
  fichiers de référence.
\item
  \textbf{Analyse univariée} : distributions, indicateurs de position et
  de dispersion, modalités dominantes.
\item
  \textbf{Décompte des valeurs manquantes} par variable, en distinguant
  les vides métier des vides d'anonymisation quand c'est possible.
\item
  \textbf{Décompte des outliers} par variable numérique, avec la méthode
  utilisée (IQR, z-score). Interprétez avant de supprimer : un solde
  négatif peut être normal.
\end{itemize}

\subsection{Bonus}\label{bonus}

\begin{itemize}
\tightlist
\item
  \textbf{Analyse bivariée} : variable explicative × churn,
  corrélations, tests statistiques.
\item
  \textbf{Web scraping} avec \textbf{Beautiful Soup} (pages statiques),
  \textbf{Playwright} (pages dynamiques) ou \textbf{Scrapy} (collecte
  structurée).
\end{itemize}

\begin{quote}
\textbf{Règles du web scraping} : collectez uniquement des informations
publiques et générales sur le secteur bancaire ou le churn. Respectez
les conditions d'utilisation et \texttt{robots.txt} des sites. Il est
interdit de chercher à relier ces données au fichier fourni ou à
identifier l'établissement ou un client. Les dates étant décalées,
aucune jointure calendaire n'a de sens.
\end{quote}

\subsection{Outils proposés}\label{outils-proposuxe9s-1}

{\def\LTcaptype{none} % do not increment counter
\begin{longtable}[]{@{}ll@{}}
\toprule\noalign{}
Catégorie & Options \\
\midrule\noalign{}
\endhead
\bottomrule\noalign{}
\endlastfoot
Manipulation & \textbf{pandas}, \textbf{polars}, \textbf{numpy} \\
Visualisation & \textbf{matplotlib}, \textbf{seaborn},
\textbf{plotly} \\
Profilage automatique & \texttt{ydata-profiling}, en local uniquement \\
\end{longtable}
}

\begin{center}\rule{0.5\linewidth}{0.5pt}\end{center}

\section{Étape 2 --- Conception du DWH et ETL (semaine du
16/11)}\label{uxe9tape-2-conception-du-dwh-et-etl-semaine-du-1611}

\subsection{Obligatoire : Talend + PostgreSQL /
pgAdmin}\label{obligatoire-talend-postgresql-pgadmin}

\begin{enumerate}
\def\labelenumi{\arabic{enumi}.}
\tightlist
\item
  \textbf{Conception} : définir le grain, la table de faits et les
  dimensions du modèle en étoile.
\item
  \textbf{Base cible} : créer la base et le schéma sous
  \textbf{PostgreSQL} avec \textbf{pgAdmin}.
\item
  \textbf{Extract} : lire le CSV et les fichiers Excel dans
  \textbf{Talend}.
\item
  \textbf{Transform} : convertir les dates \texttt{YYYYMMDD}, traiter
  les vides et les doublons, joindre les fichiers de référence, calculer
  les variables dérivées.
\item
  \textbf{Load} : charger les dimensions, puis les faits, en contrôlant
  les clés.
\item
  \textbf{Contrôle} : comparer les volumes et les indicateurs entre la
  source et l'entrepôt.
\end{enumerate}

\subsection{Bonus}\label{bonus-1}

{\def\LTcaptype{none} % do not increment counter
\begin{longtable}[]{@{}
  >{\raggedright\arraybackslash}p{(\linewidth - 2\tabcolsep) * \real{0.5000}}
  >{\raggedright\arraybackslash}p{(\linewidth - 2\tabcolsep) * \real{0.5000}}@{}}
\toprule\noalign{}
\begin{minipage}[b]{\linewidth}\raggedright
Bonus
\end{minipage} & \begin{minipage}[b]{\linewidth}\raggedright
Attendu
\end{minipage} \\
\midrule\noalign{}
\endhead
\bottomrule\noalign{}
\endlastfoot
\textbf{SSIS} & Reproduire un flux et comparer avec Talend \\
\textbf{Databricks} & Architecture médaillon : \textbf{bronze} (brut),
\textbf{silver} (nettoyé), \textbf{gold} (prêt pour l'analyse) \\
\textbf{SCD de type 2} & Historiser une dimension avec dates de validité
et indicateur de version courante \\
\textbf{Python} & Pipeline reproductible (pandas, SQLAlchemy) \\
\end{longtable}
}

\begin{quote}
\textbf{Important} : le pipeline doit être \textbf{reproductible} et
\textbf{documenté} par un schéma clair du flux. Exportez les jobs Talend
dans le dépôt, sans les fichiers de données. Si vous utilisez
Databricks, n'y chargez les données qu'avec l'accord des tuteurs.
\end{quote}

\subsection{Outils proposés}\label{outils-proposuxe9s-2}

{\def\LTcaptype{none} % do not increment counter
\begin{longtable}[]{@{}ll@{}}
\toprule\noalign{}
Catégorie & Options \\
\midrule\noalign{}
\endhead
\bottomrule\noalign{}
\endlastfoot
ETL & \textbf{Talend} (obligatoire) \\
Base de données & \textbf{PostgreSQL + pgAdmin} (obligatoire) \\
Schéma & dbdiagram.io, draw.io, Mermaid \\
\end{longtable}
}

\begin{center}\rule{0.5\linewidth}{0.5pt}\end{center}

\section{Étape 3 --- Reporting et dashboard (pour le
23/11)}\label{uxe9tape-3-reporting-et-dashboard-pour-le-2311}

\subsection{Obligatoire : Power BI +
DAX}\label{obligatoire-power-bi-dax}

\begin{itemize}
\tightlist
\item
  Connecter \textbf{Power BI Desktop} à PostgreSQL.
\item
  Définir les KPI en \textbf{DAX} : taux de churn, ancienneté moyenne,
  solde moyen, etc. Attention au grain : un taux calculé par ligne n'est
  pas un taux par compte.
\item
  Construire au minimum \textbf{3 pages} :

  \begin{enumerate}
  \def\labelenumi{\arabic{enumi}.}
  \tightlist
  \item
    \textbf{Vue d'ensemble} : KPI clés, taux de churn global, évolution.
  \item
    \textbf{Analyse du churn} : décomposition par segment (âge, secteur,
    ancienneté, produit, devise).
  \item
    \textbf{Segmentation client} : profils et valeur client.
  \end{enumerate}
\item
  Soigner la mise en page : filtres, navigation entre pages, infobulles.
\end{itemize}

\subsection{Bonus}\label{bonus-2}

\begin{itemize}
\tightlist
\item
  \textbf{Visualisation avec Databricks}.
\item
  \textbf{Interface JavaScript} (par exemple Chart.js, D3.js ou React)
  alimentée par des agrégats.
\end{itemize}

\begin{quote}
\textbf{Conseil} : la qualité du design compte, mais la justesse des
mesures compte davantage. Ne publiez pas le rapport sur le Web : un
\texttt{.pbix} contient des données. Remettez-le aux tuteurs hors
GitHub.
\end{quote}

\begin{center}\rule{0.5\linewidth}{0.5pt}\end{center}

\section{Étape 4 --- Modèles prédictifs (pour le
1/12)}\label{uxe9tape-4-moduxe8les-pruxe9dictifs-pour-le-112}

\subsection{4.1 Obligatoire : fixer un KPI et prévoir sa série
temporelle}\label{obligatoire-fixer-un-kpi-et-pruxe9voir-sa-suxe9rie-temporelle}

\begin{enumerate}
\def\labelenumi{\arabic{enumi}.}
\tightlist
\item
  \textbf{Choisir un KPI} utile au métier, par exemple le nombre
  hebdomadaire de clôtures de comptes ou le taux de clôture.
\item
  \textbf{Construire la série} au grain compte, en semaines, en ajoutant
  les semaines à zéro.
\item
  \textbf{Analyser} : tendance, saisonnalité, autocorrélation
  (ACF/PACF), stationnarité, pics.
\item
  \textbf{Traiter les pics} : cinq semaines concentrent environ la
  moitié des clôtures. Identifiez-les et justifiez votre traitement.
\item
  \textbf{Modéliser} : entraîner un \textbf{LSTM} ou un \textbf{RNN} sur
  des fenêtres glissantes, avec normalisation calculée uniquement sur
  l'entraînement.
\item
  \textbf{Comparer} à une référence simple (prévision naïve ou
  saisonnière, éventuellement ARIMA). Le LSTM doit prouver son intérêt :
  avec environ 200 points, il peut ne pas gagner, et ce résultat est
  acceptable s'il est démontré.
\item
  \textbf{Évaluer} sur la fin de la série, sans mélange temporel, avec
  MAE et RMSE. Le MAPE n'est pas défini sur les semaines à zéro.
\end{enumerate}

\subsection{4.2 Obligatoire : Machine Learning pour le
churn}\label{obligatoire-machine-learning-pour-le-churn}

\begin{enumerate}
\def\labelenumi{\arabic{enumi}.}
\tightlist
\item
  \textbf{Préparer} : fixer le grain, sélectionner les variables et
  écarter celles qui provoquent une \textbf{fuite de cible} (document 2,
  section 5.1).
\item
  \textbf{Modéliser} : comparer \textbf{au moins 3 modèles} (par exemple
  régression logistique, Random Forest, XGBoost).
\item
  \textbf{Découper sans contamination} : un même client ne doit pas se
  trouver à la fois dans l'entraînement et le test
  (\texttt{GroupShuffleSplit}).
\item
  \textbf{Évaluer} : precision, recall, F1, \textbf{ROC-AUC},
  \textbf{PR-AUC}, matrice de confusion. Pas seulement l'accuracy.
\item
  \textbf{Optimiser et interpréter} : hyperparamètres du meilleur
  modèle, importance des variables, valeurs SHAP.
\item
  \textbf{Sauvegarder} le modèle final hors du dépôt.
\end{enumerate}

\subsection{Bonus}\label{bonus-3}

\begin{itemize}
\tightlist
\item
  Traiter \textbf{plusieurs KPI}, avec un modèle par KPI ou un modèle
  multivarié.
\end{itemize}

\subsection{Outils proposés}\label{outils-proposuxe9s-3}

{\def\LTcaptype{none} % do not increment counter
\begin{longtable}[]{@{}ll@{}}
\toprule\noalign{}
Catégorie & Options \\
\midrule\noalign{}
\endhead
\bottomrule\noalign{}
\endlastfoot
Séries temporelles & \textbf{statsmodels}, \textbf{pandas} \\
Deep Learning & \textbf{TensorFlow / Keras}, \textbf{PyTorch} \\
Machine Learning & \textbf{scikit-learn}, \textbf{XGBoost},
\textbf{LightGBM} \\
Hyperparamètres & \texttt{GridSearchCV}, \texttt{RandomizedSearchCV},
\textbf{Optuna} \\
Interprétabilité & \textbf{SHAP}, importance native des arbres \\
Sauvegarde & \textbf{joblib} \\
\end{longtable}
}

\begin{center}\rule{0.5\linewidth}{0.5pt}\end{center}

\section{Étape 5 --- Déploiement (pour le
7/12)}\label{uxe9tape-5-duxe9ploiement-pour-le-712}

\subsection{Obligatoire : interface
web}\label{obligatoire-interface-web}

\begin{itemize}
\tightlist
\item
  Concevoir une interface simple avec au moins :

  \begin{itemize}
  \tightlist
  \item
    un \textbf{formulaire de prédiction} : saisie d'un profil fictif →
    score de churn ;
  \item
    l'\textbf{affichage de la prévision du KPI} ;
  \item
    des \textbf{KPI de synthèse} agrégés.
  \end{itemize}
\item
  Déployer sur une plateforme accessible et donner l'URL dans le README.
\end{itemize}

\begin{quote}
\textbf{Protection des données} : une application publique n'affiche ni
ligne client, ni identifiant, ni export. Le dépôt et l'image déployée ne
contiennent aucune donnée. Si vous présentez une liste de comptes à
risque, générez-la sur des \textbf{profils fictifs} ou placez-la
derrière une authentification, sur un serveur non public.
\end{quote}

\subsection{Obligatoire : diagrammes
UML}\label{obligatoire-diagrammes-uml}

{\def\LTcaptype{none} % do not increment counter
\begin{longtable}[]{@{}
  >{\raggedright\arraybackslash}p{(\linewidth - 2\tabcolsep) * \real{0.5000}}
  >{\raggedright\arraybackslash}p{(\linewidth - 2\tabcolsep) * \real{0.5000}}@{}}
\toprule\noalign{}
\begin{minipage}[b]{\linewidth}\raggedright
Diagramme
\end{minipage} & \begin{minipage}[b]{\linewidth}\raggedright
Attendu
\end{minipage} \\
\midrule\noalign{}
\endhead
\bottomrule\noalign{}
\endlastfoot
\textbf{Cas d'utilisation global} & Acteurs et fonctionnalités de
l'application \\
\textbf{Séquence} (un seul) & Le scénario principal, par exemple la
demande de prédiction \\
\textbf{Classes} & Structure de l'application et liens entre les
classes \\
\end{longtable}
}

\subsection{Bonus}\label{bonus-4}

\begin{itemize}
\tightlist
\item
  \textbf{DevOps} : tests automatisés (\texttt{pytest}) et intégration
  continue (GitHub Actions).
\item
  \textbf{Application mobile} : Flutter, React Native ou PWA, avec les
  mêmes règles de protection des données.
\end{itemize}

\subsection{Outils proposés}\label{outils-proposuxe9s-4}

{\def\LTcaptype{none} % do not increment counter
\begin{longtable}[]{@{}ll@{}}
\toprule\noalign{}
Catégorie & Options \\
\midrule\noalign{}
\endhead
\bottomrule\noalign{}
\endlastfoot
Application data & \textbf{Streamlit}, \textbf{Gradio}, \textbf{Dash} \\
Front + API & \textbf{FastAPI} ou Flask avec React, Vue ou Bootstrap \\
Hébergement & Streamlit Community Cloud, Hugging Face Spaces, Render \\
UML & \textbf{draw.io}, PlantUML, Mermaid, StarUML \\
\end{longtable}
}

\begin{center}\rule{0.5\linewidth}{0.5pt}\end{center}

\section{Étape 6 --- Rapport et
présentation}\label{uxe9tape-6-rapport-et-pruxe9sentation}

\subsection{Objectifs}\label{objectifs-1}

\begin{itemize}
\tightlist
\item
  Documenter votre démarche, vos choix et vos résultats.
\item
  Convaincre le jury en 20 minutes que votre solution répond au besoin.
\end{itemize}

\subsection{Structure recommandée du rapport (25--40
pages)}\label{structure-recommanduxe9e-du-rapport-2540-pages}

\begin{enumerate}
\def\labelenumi{\arabic{enumi}.}
\tightlist
\item
  \textbf{Page de garde et résumé exécutif} (1 page).
\item
  \textbf{Introduction et contexte} (2 pages).
\item
  \textbf{Compréhension des données} (3--5 pages : univarié, manquants,
  outliers).
\item
  \textbf{Architecture de la solution} (2--3 pages, avec schémas).
\item
  \textbf{DWH et ETL} (3--5 pages : modèle en étoile, Talend,
  PostgreSQL).
\item
  \textbf{Reporting Power BI et DAX} (3--5 pages, captures sans ligne
  client).
\item
  \textbf{Modèles prédictifs} (5--8 pages : KPI, série temporelle,
  LSTM/RNN, churn, contrôle de la fuite de cible).
\item
  \textbf{Interface web, déploiement et UML} (2--3 pages, captures,
  diagrammes, lien démo).
\item
  \textbf{Limites et perspectives} (1--2 pages).
\item
  \textbf{Conclusion} (1 page).
\item
  \textbf{Annexes} : références, code clés, glossaire.
\end{enumerate}

\subsection{Outils proposés}\label{outils-proposuxe9s-5}

{\def\LTcaptype{none} % do not increment counter
\begin{longtable}[]{@{}
  >{\raggedright\arraybackslash}p{(\linewidth - 2\tabcolsep) * \real{0.5000}}
  >{\raggedright\arraybackslash}p{(\linewidth - 2\tabcolsep) * \real{0.5000}}@{}}
\toprule\noalign{}
\begin{minipage}[b]{\linewidth}\raggedright
Catégorie
\end{minipage} & \begin{minipage}[b]{\linewidth}\raggedright
Options
\end{minipage} \\
\midrule\noalign{}
\endhead
\bottomrule\noalign{}
\endlastfoot
Rédaction du rapport & \textbf{LaTeX} (Overleaf), \textbf{Microsoft
Word}, \textbf{Google Docs}, \textbf{Typst}, Markdown + Pandoc \\
Diagrammes & \textbf{draw.io}, \textbf{Lucidchart}, \textbf{Mermaid},
\textbf{Excalidraw}, dbdiagram.io \\
Captures d'écran & ShareX, Snipping Tool, macOS Screenshot \\
Slides & \textbf{PowerPoint}, \textbf{Google Slides}, \textbf{Canva},
\textbf{Beamer} (LaTeX), Marp (Markdown) \\
\end{longtable}
}

\subsection{Conseils pour la présentation
orale}\label{conseils-pour-la-pruxe9sentation-orale}

\begin{itemize}
\tightlist
\item
  \textbf{Temps} : 20 minutes de présentation + 10 minutes de questions.
  Ne dépassez pas.
\item
  \textbf{Rythme} : \textasciitilde1 slide par minute.
\item
  \textbf{Démo} : prévoyez \textbf{au moins 5 minutes} de démonstration
  live de votre application web. Préparez un plan B (vidéo de secours)
  en cas de problème de connexion.
\item
  \textbf{Répartition} : chaque membre de l'équipe doit prendre la
  parole.
\item
  \textbf{Public} : présentez à un public à la fois technique et métier
  --- expliquez les choix techniques de manière accessible.
\end{itemize}

\begin{center}\rule{0.5\linewidth}{0.5pt}\end{center}

\section{Récapitulatif --- Stack du
projet}\label{ruxe9capitulatif-stack-du-projet}

{\def\LTcaptype{none} % do not increment counter
\begin{longtable}[]{@{}
  >{\raggedright\arraybackslash}p{(\linewidth - 4\tabcolsep) * \real{0.3333}}
  >{\raggedright\arraybackslash}p{(\linewidth - 4\tabcolsep) * \real{0.3333}}
  >{\raggedright\arraybackslash}p{(\linewidth - 4\tabcolsep) * \real{0.3333}}@{}}
\toprule\noalign{}
\begin{minipage}[b]{\linewidth}\raggedright
Composante
\end{minipage} & \begin{minipage}[b]{\linewidth}\raggedright
Obligatoire
\end{minipage} & \begin{minipage}[b]{\linewidth}\raggedright
Bonus possibles
\end{minipage} \\
\midrule\noalign{}
\endhead
\bottomrule\noalign{}
\endlastfoot
Compréhension des données & Python, pandas, Jupyter & Scraping, analyse
bivariée \\
ETL & \textbf{Talend} & SSIS, Databricks, Python \\
Data Warehouse & \textbf{PostgreSQL + pgAdmin} & Médaillon, SCD de type
2 \\
Reporting & \textbf{Power BI + DAX} & Databricks, JavaScript \\
Série temporelle & \textbf{LSTM / RNN} & Plusieurs KPI \\
Churn & Machine Learning (scikit-learn, XGBoost) & Plusieurs KPI \\
Interface & Application web + \textbf{UML} & Mobile, tests DevOps \\
Versioning & \textbf{Git + GitHub} & GitHub Actions \\
\end{longtable}
}

\end{document}
